\subsection{Trazas con cotas independientes de la dimensión}\label{sec:B5T04}
Para $t(X)=\sum_{ab}|X_{ab}|$,
\[
 t(AXB)\le\|A\|_1t(X)\|B\|_\infty.
\]
Sean $D=S_*+B_*$ y $H=I+K_*$ simétricas, con
\[
 \|S_*\|_\infty=O(k^{-1}),\ t(S_*)=O(1),\quad
 \|B_*\|_{\max}=O(k^{-2}),\ t(B_*)=O(1),
\]
\[
 \|B_*\|_\infty=O(k^{-1}),\quad
 \|K_*\|_{\max}=O(k^{-1}),\quad\|H\|_\infty=O(1).
\]
Entonces $t(S_*^2)=O(k^{-1})$, y
\[
 |\operatorname{tr}K_*S_*^2|=O(k^{-2}),\quad
 \|K_*S_*\|_{\max}=O(k^{-2}),\quad t(K_*S_*)=O(1).
\]
Luego $\operatorname{tr}(HS_*HS_*)=\operatorname{tr}S_*^2+O(k^{-2})$.
Los términos con $B_*$ se acotan por su máximo multiplicado por
$t(HS_*H)$ o $t(HB_*H)$. Se obtiene
\[
 \boxed{\operatorname{tr}(HDHD)=\operatorname{tr}S_*^2+O(k^{-2}).}
\]
Además $\|HD\|_{op}=O(k^{-1})$ y $\|HD\|_{HS}^2=O(k^{-1})$.
La desigualdad $|\operatorname{tr}T^r|\le\|T\|_{op}^{r-2}\|T\|_{HS}^2$
da un resto cúbico $O(k^{-2})$ en $\log\det(I+HD)$.

\paragraph{Aplicación al cambio de soporte.}
En el defecto $D_*$ con gauge exacto de la Sección~\ref{sec:B4T08}, $S_*$ contiene
las entradas $-p_i$ de las casillas $i\in F\setminus\{c\}$.
Cada una de ellas toca una línea retirada y tres retenidas, por lo que aporta seis
entradas orientadas. Las intersecciones únicas impiden solapamientos.
Así
\[
 \operatorname{tr}S_*^2=6\sum_Fp_i^2+O(k^{-2}).
\]
Los restantes términos proceden de $y$, del término de Schur de rango cuatro y del cambio del
proyector RC, y satisfacen las cotas de $B_*$ por las Secciones~\ref{sec:B5T02} y~\ref{sec:B5T03}.
El inverso $H_*$ es $(\Gamma+\Pi)_{II}$ y satisface las cotas de $H$.
Por tanto se conserva el coeficiente seis, no sólo su orden.

\paragraph{Aplicación al prefactor.}
Sea $E=S-S_0=A^T\operatorname{diag}(3p^2+12p^3+O(p^4))A$.
Se tiene $t(E)=O(1)$, $\|E\|_\infty=O(k^{-1})$,
$E_{ab}=O(k^{-2})$ para $a\ne b$, y
\[
 E_{aa}=3m_{2,a}+r_a,\quad\sum_a|r_a|=O(k^{-1}),
 \quad\max_a|r_a|=O(k^{-2}).
\]
Al separar la diagonal de la parte fuera de diagonal resulta
\[
 \operatorname{tr}(\Gamma E\Gamma E)
 =9\sum_a m_{2,a}^2+O(k^{-2}),\qquad
 \text{resto cúbico}=O(k^{-2}).
\]
Para el cambio padre--hijo se extiende $m'_2$ por cero en las líneas retiradas.
Contando cada casilla a lo sumo cuatro veces,
\[
 \sum_a|\Delta m_{2,a}|\le
 4\sum_Fp_i^2+4\sum_{E'}(2p_i|y_i|+y_i^2)+\sum_{a\in J}m_{2,a}
 =O(k^{-1}).
\]
Como $\max m_2,\max m'_2=O(k^{-1})$, resulta
$\Delta\sum_a m_{2,a}^2=O(k^{-2})$.

\paragraph{Resolvente por entradas.}
Con $\Delta_0=S'_{0}-S_{0,II}$, de diagonal cero y norma infinito $O(k^{-1})$,
los complementos de Schur y los gauges exactos dan
\[
 \Gamma'-\Gamma_{II}=-\Delta_0+O_{\max}(k^{-2}).
\]
Sustituir las inversas por $I$ en el término lineal cuesta
$\|H-I\|_{\max}\|\Delta_0\|_1\|H'\|_1=O(k^{-2})$;
los términos cuadráticos tienen ese mismo orden por normas de filas.
Entre líneas de una casilla superviviente $i$ se tiene $\Delta_{0,ab}=y_i$;
no interviene ninguna casilla retirada. Por ello
\[
 \kappa'_i-\kappa_i=O(k^{-2}),\quad
 \kappa_i=a_i^T\Gamma a_i-4,
 \quad\Delta\sum_i p_i^2\kappa_i=O(k^{-2}).
\]
También $\Delta\sum p_i^3=O(k^{-2})$ por $\sum_Fp=O(1)$ y $y/p=O(k^{-1})$.
Estas son las sustituciones de las Secciones~\ref{sec:B4T07} y~\ref{sec:B4T08}; nunca se multiplica
un error uniforme por el número de diagonales.

\paragraph{Alcance de las cotas de tamaño.}
El resto de Newton de orden relativo $k^{-2}$ tiene masa $O(k^{-1})$.
Las trazas lineales de orden $k^{-1}$ se conservan en la cancelación de la Sección~\ref{sec:B5T05};
no se desechan por estas cotas de tamaño.
